\section{Current Situation}

\subsection{Problem Statement}

The causal chain described in the “Introduction to the Topic” is not an abstract story: it has produced the developing world’s most acute debt pressures in a generation. The combined external debt of low- and middle-income countries reached an all-time high of US\$8.9 trillion in 2024, of which a record US\$1.2 trillion is owed by the 78 poorest countries eligible to borrow from IDA.\footnote{World Bank. (2025b, December 3). Developing countries’ debt outflows hit 50-year high during 2022-2024. \url{https://www.worldbank.org/en/news/press-release/2025/12/03/developing-countries-debt-outflows-hit-50-year-high-during-2022-2024}} The proportion of IDA-eligible countries assessed as being at high risk of, or already in, debt distress has risen sharply, from roughly 30\% in 2012 to approximately 58\% by 2021.\footnote{Independent Evaluation Group. (2023). How the Low-Income Country Debt Sustainability Framework informs World Bank corporate and operational decisions (Chapter 7). World Bank Group. \url{https://ieg.worldbankgroup.org/evaluations/world-banks-role-and-use-low-income-country-debt-sustainability-framework}} Well over half of the world’s poorest countries are now in or near a debt trap, leaving less money for development spending.

What makes this crisis particularly hard to solve is that no single institution can resolve it alone. In past crises, debt was largely owed to a coordinated group of Western governments (the Paris Club) and could be restructured at a single table. Today, debt is owed to a fragmented creditor landscape, with China now being the largest single bilateral lender to many developing countries.\footnote{Lowy Institute. (2025, May). Peak repayment: China’s global lending. \url{https://interactives.lowyinstitute.org/features/peak-repayment-china-global-lending/}} Additionally, private bondholders, a class with no collective body to negotiate with, are the fastest-growing class of creditors, even in the poorest IDA economies.\footnote{World Bank. (2026, February 17). Beyond the narratives: What the IDR 2025 data reveal about debt in low- and middle-income countries. \url{https://blogs.worldbank.org/en/opendata/beyond-the-narratives--what-the-idr-2025-data-reveal-about-debt-}} This fragmentation, combined with higher refinancing rates, has turned a large debt stock into a complex crisis to solve.

\subsection{Scope: What This Committee Can and Cannot Do}

Because sovereign debt is a crowded field of institutions, delegates must be disciplined about staying inside the World Bank’s mandate. The Bank is not the body that restructures sovereign debt, declares countries bankrupt, provides emergency balance-of-payments bailouts, or sets monetary and exchange-rate policy. Those functions belong to the IMF, to creditor groupings such as the Paris Club and the G20, and ultimately to debtor governments and their creditors negotiating directly.\footnote{International Monetary Fund. (2022). The IMF and the World Bank (factsheet). \url{https://www.imf.org/en/about/factsheets/sheets/2022/imf-world-bank-new}} The Bank nonetheless sits inside that process, as a creditor itself, as the provider of the debt data on which restructuring depends, and as the institution that protects development finance during crises.\footnote{Independent Evaluation Group. (2023). How the Low-Income Country Debt Sustainability Framework informs World Bank corporate and operational decisions (Chapter 7). World Bank Group. \url{https://ieg.worldbankgroup.org/evaluations/world-banks-role-and-use-low-income-country-debt-sustainability-framework}}

What the Bank can do is the substance of this committee, and it is broader than may first appear. When developing solutions, delegates should keep in mind that the Bank’s instruments only work through its own levers, or by indirectly shaping the choices other actors make. It can never force an outcome on a country or institution that is not its own. Countries listen to the Bank because they depend on its financing, policy advice, and technical expertise. At the same time, how it treats a country’s debt directly shapes how international markets evaluate that country’s creditworthiness. This gives the Bank real influence, and with it, real responsibility for how that influence is used. In practice, this means it can adjust the terms and conditions attached to its lending, introduce new standards for creditors and borrowers, or build cooperation among actors toward common responses. A more detailed description of how concrete measures can unfold is provided in “The Bank’s Toolkit” section below. The art of this committee is to craft solutions that tackle the underlying issues of debt distress while staying within the Bank’s powers.

\subsection{The Bank’s Toolkit}

This section gives three concise examples of specific tools the World Bank can use. It is by no means exhaustive, and delegates may find others that serve as useful building blocks for their own solutions. A good solution names the problem it addresses, sets out a workable mechanism for how it would function, ideally builds on or enhances something that already exists, and has a real impact on the issue the committee is debating.

Firstly, the World Bank, together with the IMF, prepares Debt Sustainability Analyses (DSAs). A DSA is essentially a checkup of a country’s financial health that evaluates whether it can manage its debt obligations sustainably. Although the DSA does not dictate the outcome of a case, it serves as the basis for discussion at the creditor table and is commonly accepted as a shared framework for assessing a nation’s debt situation. It also stress-tests the country’s economy against shocks, and places it in one of four categories: low risk, moderate risk, high risk, or already in debt distress.\footnote{International Monetary Fund. (2025). IMF-World Bank Debt Sustainability Framework for low-income countries (factsheet). \url{https://www.imf.org/en/about/factsheets/sheets/2023/imf-world-bank-debt-sustainability-framework-for-low-income-countries}} This rating is important because it determines the mix of IDA grants and loans the country may receive.\footnote{International Development Association. (2023). Debt. World Bank. \url{https://ida.worldbank.org/en/financing/debt}}

Secondly, the Sustainable Development Finance Policy (SDFP) gives the Bank a lever to drive behavioural change in IDA-eligible countries, e.g., by incentivising them to implement policies that promote more sustainable debt management. The SDFP operates through two linked pillars: The Debt Sustainability Enhancement Program (DSEP) and the Program of Creditor Outreach (PCO), which are beyond the scope of this section but delegates may find worth researching independently. The World Bank can use the SDFP to expose hidden debt, modernise tax collection or prevent a country from taking on harmful deals with other lenders. In cases in which countries fail to meet these targets, the Bank can withhold its funding. In summary, this tool supports the alignment of global lending practices to ensure that countries commit to policies that help them overcome a crisis or prevent it from happening in the first place.\footnote{World Bank. (2025c, May 6). Sustainable Development Finance Policy of the International Development Association: FY25 implementation update and policy enhancements. \url{https://documents.worldbank.org/en/publication/documents-reports/documentdetail/099051625184539973}}

Thirdly, debt-for-development and debt-for-nature swaps allow countries to exchange expensive debt for cheaper debt, with the condition that the money saved goes to an agreed purpose. The difference between these swaps and the SDFP is that, unlike the SDFP, the swaps are financial transactions that actively refinance old debt into cheaper debt. In December 2024, Côte d’Ivoire was the first country to receive World-Bank-backed debt-for-development swaps. The World Bank provided a €240 million financial guarantee that convinced commercial lenders to offer better terms, which allowed the government to replace close to €400 million of expensive debt with a cheaper loan. This deal freed up €330 million over five years, which is now being allocated to the development of over 30 schools for 30,000 students.\footnote{World Bank Group. (2025, February 27). Innovative finance for education: How Côte d’Ivoire’s debt swap is creating new opportunities (MIGA case study). \url{https://www.miga.org/case-study/innovative-finance-education-how-cote-divoires-debt-swap-creating-new-opportunities}} Despite this success, a limitation is that swaps require a country to be solvent, meaning that they are not applicable to countries that have already defaulted.\footnote{World Bank. (2024d, December 5). How are countries trading debt for development? A closer look at debt swaps. \url{https://www.worldbank.org/en/topic/debt/brief/how-are-countries-trading-debt-for-development}}

Other instruments worth researching include:

\begin{itemize}
  \item Development Policy Financing (DPF), Program-for-Results Lending (PforR), and Investment Project Financing (IPF): how the Bank disburses financing
  \item The IDA Debt Reduction Facility (DRF): buying back commercial debt at a discount
  \item Climate Resilient Debt Clauses (CRDCs): pausing payments after a disaster
  \item Catastrophe Deferred Drawdown Options (Cat DDOs): contingent financing after a disaster
  \item The Debtor Reporting System (DRS) and the Debt Management Facility (DMF): debt data and capacity
  \item Global Sovereign Debt Roundtable (GSDR): World Bank, IMF, and G20-led dialogue on creditor-coordination bottlenecks
  \item Blended and local-currency finance: reducing financing costs and foreign exchange risk
  \item Domestic Resource Mobilisation (DRM) support: strengthening tax collection
\end{itemize}
